\documentclass[a4paper,11pt]{article}

\usepackage[T1]{fontenc}
\usepackage{lmodern}
\usepackage[french]{babel}
\usepackage[top=2.2cm, bottom=2.5cm, left=2.5cm, right=2.5cm]{geometry}
\usepackage{amsmath,amssymb}
\usepackage[dvipsnames]{xcolor}
\usepackage{fancyhdr}
\usepackage{lastpage}
\usepackage{enumitem}
\usepackage{hyperref}
\hypersetup{colorlinks=true, linkcolor=MidnightBlue, urlcolor=MidnightBlue}

\newcommand*\dif{\mathop{}\!\mathrm{d}}
\newcommand*\diver{\mathop{}\!\mathrm{div}}

\pagestyle{fancy}
\renewcommand\headrulewidth{0.6pt}
\fancyhead[L]{Relecture des exercices de colle}
\fancyhead[R]{Matthieu Santin}
\fancyfoot[C]{\thepage/\pageref{LastPage}}

% Marqueurs de relecture
\newcommand{\exo}[1]{\vspace{0.7\baselineskip}\noindent\textbf{\large #1}\par\nopagebreak\vspace{0.2\baselineskip}}
\newcommand{\erreur}{\textcolor{BrickRed}{\textbf{Erreur}}}
\newcommand{\suggestion}{\textcolor{MidnightBlue}{\textbf{Suggestion}}}
\newcommand{\verifie}{\textcolor{ForestGreen}{\textbf{Vérifié}}}
\newcommand{\rien}{\textcolor{ForestGreen}{\textbf{Rien à corriger.}}}
\newcommand{\traite}{\textcolor{ForestGreen}{\textbf{Traité}}}

\setlist[itemize]{leftmargin=1.4em, itemsep=0.25\baselineskip, topsep=0.3\baselineskip, label=\textbullet}

\title{\vspace{-1.5cm}Relecture des exercices de colle\\[0.3em]\large Recueil \texttt{main.pdf}}
\author{}
\date{Septembre 2026}

\begin{document}

\maketitle
\thispagestyle{fancy}

\noindent\textbf{Objet.} Relecture exercice par exercice du recueil \texttt{exercices\_colle/main.pdf} :
vérification des calculs et des applications numériques, contrôle des corrigés contre les schémas,
cohérence des renvois aux figures, et suggestions de fond ou de rédaction.

\medskip
\noindent\textbf{Méthode.} Chaque résultat encadré et chaque valeur numérique sont recalculés
indépendamment. Les schémas sont relus pour vérifier que la topologie décrite dans le corrigé est
bien celle qui est dessinée, et les renvois aux couleurs des courbes sont contrôlés sur les figures.

\medskip
\noindent\textbf{Conventions.} \erreur{} signale ce qui est faux et doit être repris.
\suggestion{} signale une amélioration possible, sans urgence. \verifie{} récapitule ce qui a été
recalculé et confirmé, pour mémoire.

\medskip
\noindent\textbf{Suivi.} \traite{} en tête de chapitre signifie que les corrections ont
déjà été reportées dans \texttt{main.pdf}. Le détail reste consigné ici pour mémoire.

\medskip
\noindent\textbf{Périmètre.} Les sections 1.5 à 1.8 du chapitre \'Electronique ont été rédigées au
cours de ce travail de relecture et ne sont pas reprises ici.

\section{\'Electronique}

Les exercices 1.2, 1.3 et 1.4 forment une série cohérente : même structure de questions, et surtout
\textbf{la même pulsation propre} $\omega_0 = 6{,}29\times10^{6}$~rad$\cdot$s$^{-1}$ dans les trois cas,
obtenue avec des composants différents. C'est un choix heureux, qui permet de les donner en
comparaison et de faire sentir qu'à fréquence propre identique trois filtres agissent très
différemment sur le même créneau. Autant l'assumer explicitement dans une phrase d'introduction au
groupe.

\exo{1.1 Générateur de triangles (Mines-Ponts PSI 2022)}

\rien{}

\verifie{} : les régimes de fonctionnement des trois ALI contre le schéma (rétroaction sur $-$ pour
$(A1)$ et $(A2)$, sur $+$ seulement pour $(A3)$), la relation $v_1=-\frac{R_2}{R_1}v_e$, la relation
d'intégrateur, le sens horaire du cycle d'hystérésis, la durée $t_b=\frac{2R_1R_3RC}{R_2(R_3+R_4)}$,
la période $T=\frac{4R_1R_3RC}{R_2(R_3+R_4)}$ et l'application numérique ($v_b=7{,}5$~V,
$T=200~\mu$s, $f=5$~kHz).

\begin{itemize}
	\item \suggestion{} Cinq questions dont deux tracés à main levée : c'est long pour une colle de
	vingt minutes. Si l'exercice est donné en entier, autant fournir le résultat de la question 2 pour
	que l'élève démarre à la question 3.
	\item \suggestion{} Le titre mentionne « Mines-Ponts PSI 2022 » sans préciser l'épreuve. La base
	\texttt{annales/base/} permet de retrouver le sujet ; indiquer l'épreuve et la partie faciliterait
	le retour à l'énoncé d'origine.
\end{itemize}

\exo{1.2 Effet d'un filtre passif $RLC$ parallèle}

\erreur{} \textbf{question 1.} Le corrigé écrit :

\begin{quote}
« dans les deux cas l'impédance de la branche série $L\parallel C$ tend vers zéro, \textit{aucun
courant ne circule dans $R$} et $V_s = V_e$ »
\end{quote}

Les deux affirmations sont incompatibles. Si l'impédance de la branche série tend vers zéro, celle-ci
se comporte comme un fil : la totalité de $V_e$ se retrouve aux bornes de $R$, et le courant $V_e/R$
y circule bel et bien. La conclusion $V_s=V_e$ est juste, c'est l'argument qui ne l'est pas. Il faut
dire que la branche série ne présente \textit{aucune chute de tension}.

À noter que la tournure est correctement employée dans les exercices 1.3 et 1.4, où le courant est
réellement nul parce que c'est la branche \textit{série} qui est ouverte. C'est probablement un
report machinal d'un exercice à l'autre.

\medskip
\verifie{} : $\omega_0=1/\sqrt{LC}=6{,}29\times10^{6}$~rad$\cdot$s$^{-1}$, $Q=R\sqrt{C/L}=15{,}9$,
les quatre lignes du tableau de la question 4 (modules et phases recalculés harmonique par
harmonique), l'équation différentielle, les deux conditions initiales, la solution
$V_s=E_0\left[1-\frac{2\mu}{\Omega}e^{-\mu t}\sin\Omega t\right]$ et l'identité $2\mu/\Omega\simeq1/Q$.
Les renvois aux couleurs (entrée orange, sortie rouge en pointillé) correspondent bien à la figure.

\exo{1.3 Effet d'un filtre passif $RC$}

\rien{}

\verifie{} : $\underline{H}=\frac{jRC\omega}{1+3jRC\omega-(RC\omega)^2}$, $H_0=Q=1/3$,
$\omega_0=1/RC=6{,}29\times10^{6}$~rad$\cdot$s$^{-1}$, le gain de $-9{,}5$~dB à la résonance, les
quatre lignes du tableau, la décroissance des amplitudes de sortie en $1/n^2$, l'équation
différentielle, les conditions initiales, la solution $V_s=\frac{E_0}{\mu'RC}\sinh(\mu't)e^{-\mu t}$
avec $\mu=\frac{3}{2RC}$ et $\mu'=\frac{\sqrt5}{2RC}$, et le temps de décroissance
$1/(\mu-\mu')=0{,}42~\mu$s. Les couleurs (entrée noire, sortie verte en pointillé) et le facteur de
dilatation $10$ de la figure sont cohérents avec l'amplitude annoncée de $0{,}14$~V.

\begin{itemize}
	\item \suggestion{} $C=1{,}59\times10^{-2}$~nF se lit mal. $15{,}9$~pF serait plus naturel et
	donnerait tout de suite le bon ordre de grandeur.
	\item \suggestion{} La remarque « le facteur de qualité $Q=1/3$ est indépendant de $R$ et $C$ »
	mérite la précision « à composants appariés » : elle suppose les deux résistances égales entre
	elles et les deux condensateurs égaux entre eux. La figure le montre, le texte ne le dit pas.
	\item \suggestion{} Le pont de Wien étant nommé, un mot de plus situerait l'exercice : c'est le
	réseau sélectif de l'oscillateur à pont de Wien, dont le gain $1/3$ à la résonance impose un
	amplificateur de gain $3$. Cela fait le lien avec l'oscillateur quasi-sinusoïdal du \S~1.8.
\end{itemize}

\exo{1.4 Effet d'un filtre passif $RLC$ série}

\rien{}

\verifie{} : $\omega_0=6{,}29\times10^{6}$~rad$\cdot$s$^{-1}$, $Q=\frac{1}{R}\sqrt{L/C}=0{,}629$, les
quatre lignes du tableau, la largeur de bande $\Delta\omega=\omega_0/Q>\omega_0$, le caractère
pseudo-périodique de la réponse indicielle, $\Omega=3{,}81\times10^{6}$~rad$\cdot$s$^{-1}$, le
maximum de la réponse à $0{,}68\,E_0$ (recalculé par annulation de la dérivée), la constante de temps
$\tau=1/\mu=0{,}2~\mu$s, et la cohérence entre les pics de la question 4 ($0{,}75$~V) et le double de
la réponse indicielle ($1{,}35\,E_0$). Les couleurs (entrée noire, sortie bleue en pointillé)
correspondent à la figure.

\begin{itemize}
	\item \suggestion{} La discussion de la question 4 est la plus riche des trois exercices, parce
	qu'elle explique pourquoi l'approximation du dérivateur est \textit{grossière} ici. Elle gagnerait
	à être signalée comme le point d'arrivée de la série : c'est là qu'on voit qu'un filtre peu
	sélectif ne fait ni tout à fait un dérivateur ni tout à fait rien.
\end{itemize}

\section{Puissance électrique en régime sinusoïdal}

\erreur{} \textbf{titre du chapitre.} \texttt{régime sinuo\"idal} : il manque un \texttt{s}.
La coquille apparaît dans le titre du chapitre et dans la table des matières de \texttt{main.pdf}.

\medskip
Le chapitre ne contient qu'un seul exercice. Soit il mérite d'être étoffé, soit il gagnerait à être
rattaché au chapitre \'Electronique, quitte à y former une section à part.

\exo{2.1 Puissance absorbée par une installation}

\erreur{} \textbf{question 1, relation tension-courant inversée.} Le corrigé écrit
$\underline{i_1}=(R_1+jL_1\omega)\,\underline{u}$. C'est l'inverse : l'impédance multiplie le courant,
pas la tension, et il faut lire $\underline{u}=(R_1+jL\omega)\,\underline{i_1}$. La formule qui suit,
$L=\frac{1}{\omega}\sqrt{U^2/I_1^2-R_1^2}$, est quant à elle correcte, c'est bien
$|\underline{Z}|=U/I_1$ qui est utilisé. Seule la ligne intermédiaire est fausse, mais c'est
exactement le genre d'inversion qu'un élève recopie.

\medskip
\erreur{} \textbf{question 1, déphasage sans signe.} $i_1(t)=\sqrt{2}I_1\cos(\omega t\,\varphi_1)$ :
il manque l'opérateur. La charge étant inductive, le courant est en retard et il faut
$\cos(\omega t-\varphi_1)$.

\medskip
\erreur{} \textbf{question 2, $R_2$ n'est jamais calculée.} L'énoncé demande explicitement
« calculer les valeurs efficaces $I_2$ et $I$, \textit{la résistance $R_2$}, la puissance $P$ et le
facteur de puissance ». Le corrigé traite les quatre autres grandeurs et oublie celle-là. Il
manque $R_2=U^2/P_2=48{,}4~\Omega$.

\medskip
\erreur{} \textbf{notation de l'inductance, trois écritures pour une même grandeur.} Le préambule
de l'énoncé parle de $(R_1,L_1)$, la question 1 demande « l'inductance $L$ », la figure porte $L$ sur
les trois schémas, et le corrigé utilise $L_1$. Comme la figure est la plus difficile à modifier,
$L$ semble le choix à retenir partout.

\medskip
\verifie{} : $R_1=P_1/I_1^2=6{,}04~\Omega$, $\cos\varphi_1=P_1/(UI_1)=0{,}4995$ donc bien
$\varphi_1=\pi/3$, $L=33{,}3$~mH, $I_2=P_2/U=4{,}55$~A, $I=20{,}9$~A par somme vectorielle,
$\cos\varphi_2=0{,}654$ (recoupé par $P/(UI)$, qui donne la même valeur), $P=3{,}0$~kW, et
$C=I_1\sin\varphi_1/(U\omega)=2{,}28\times10^{-4}$~F. Le diagramme de Fresnel est cohérent avec le
texte, et il contient bien $I_3$ : le renvoi « cf schéma ci-dessus » de la question 3 est donc
légitime.

\begin{itemize}
	\item \suggestion{} \textbf{Le chapitre s'appelle « puissance en régime sinusoïdal » et seule la
	puissance active y apparaît.} Introduire la puissance réactive $Q=UI\sin\varphi$ et la puissance
	apparente $S=UI$ rendrait la question 3 immédiate : le condensateur doit fournir
	$Q_C=Q_1=P_1\tan\varphi_1=3{,}46$~kVAR, d'où $C=\frac{P_1\tan\varphi_1}{U^2\omega}$, soit la même
	valeur sans passer par le diagramme de Fresnel. C'est la méthode standard en électrotechnique, et
	elle se généralise à un nombre quelconque de récepteurs.
	\item \suggestion{} \textbf{$R_1$ n'est pas une résistance dissipative.} Le corrigé écrit que « la
	puissance est, en moyenne, uniquement dissipée dans la résistance ». C'est vrai du \textit{modèle
	équivalent}, mais l'installation est faite de machines tournantes : l'essentiel de ces $2$~kW part
	en travail mécanique, pas en chaleur. La confusion entre le schéma équivalent et la réalité
	physique est un classique, autant la désamorcer en une phrase.
	\item \suggestion{} La condition imposée par le distributeur mérite d'être précisée. En France on
	n'impose pas $\cos\varphi=1$ mais on facture l'énergie réactive au-delà de $\tan\varphi=0{,}4$.
	L'explication donnée, les pertes en ligne, est la bonne ; la règle de facturation la rendrait
	concrète.
	\item \suggestion{} Un ordre de grandeur sur le résultat final aiderait : $C=228~\mu$F sous
	$220$~V, c'est une batterie de condensateurs de $3{,}5$~kVAR, encombrante et coûteuse pour une
	installation de $3$~kW. C'est précisément pourquoi on se contente en pratique de relever
	$\cos\varphi$ à $0{,}93$ au lieu de $1$.
	\item \suggestion{} Le diagramme de Fresnel est introduit à la question 2 alors qu'il porte déjà
	$I_3$, qui ne sert qu'à la question 3. Soit le signaler d'un mot, soit le scinder en deux figures.
	\item \suggestion{} \textbf{Unités.} Aucune espace insécable avant les unités dans tout l'exercice
	(\texttt{\$U=220\$V}, \texttt{\$f=50\$Hz}, \texttt{\$P\_1=2,0\$kW}), là où le reste de
	\texttt{main.pdf} écrit \texttt{\$220\$\textasciitilde V}. À harmoniser.
	\item \suggestion{} Coquilles : « distribiteur », « Pourquoi le distributeur impose cette
	condition ? » (il manque l'inversion : « impose-t-il »), « les deux puissances dissipée » (accord),
	et $P$ écrit pour $P_1$ dans la première relation du corrigé.
\end{itemize}

\section{Conversion électronique de puissance}

\traite{} \textbf{le 14 septembre 2026.} Les quatre erreurs et les deux relâchements de calcul
relevés ci-dessous sont corrigés dans \texttt{main.pdf}, ainsi que les suggestions qui ne modifiaient
pas la nature des questions : hypothèse de régime permanent et convention d'ondulation crête à crête
ajoutées à la liste des hypothèses, notation $\langle I_e\rangle_{\max}$, unités et coquilles,
vérification numérique des ondulations obtenues, et renommage de
\texttt{conversion\_puissance\_electronique1\_2.pdf}.

\medskip
Deux points ont été traités autrement que suggéré. La question 7 continue de fixer $V_e=36$~V comme
dans l'énoncé d'origine, mais le mot « maximale » a été retiré et le cas le plus défavorable
($C_2>53~\mu$F) fait désormais l'objet d'une remarque pour le colleur. Les notations $V_S$, $I_S$ et
$V_E$ des figures n'ont pas été touchées, faute de source modifiable ; seul le texte a été harmonisé.

\medskip
Le détail des erreurs est conservé ci-dessous.

\exo{3.1 Hacheur, convertisseur SEPIC}

Bon exercice, bien ancré dans une situation réelle, et dont le fil conducteur (une tension d'entrée
qui traverse la tension de sortie) justifie pleinement le choix de la topologie SEPIC. Le schéma est
juste et complet. Le corrigé, en revanche, porte plusieurs confusions d'indices.

\medskip
\erreur{} \textbf{question 3, $L_1$ écrit pour $L_2$.} Après avoir correctement établi
$V_{L_2}+V_{C_1}=0$ pendant la première phase, le corrigé conclut
$V_{L_1}(0<t<\alpha T)=-V_e$. C'est $V_{L_2}$ qu'il faut lire. La question porte d'ailleurs sur
$V_{L_2}$, et le chronogramme donné juste après est bien celui de $V_{L_2}$. La contradiction est
flagrante avec la question 6, qui établit, correctement, que $V_{L_1}=+V_e$ sur ce même intervalle.

\medskip
\erreur{} \textbf{question 6, expression inversée.} Le corrigé écrit
\begin{align*}
	L_1>10\,\frac{V_s}{fI_s}\cdot\frac{\alpha^2}{1-\alpha}
	\qquad\text{alors qu'il faut}\qquad
	L_1>10\,\frac{V_s}{fI_s}\cdot\frac{(1-\alpha)^2}{\alpha}
\end{align*}
en substituant $V_e=V_s(1-\alpha)/\alpha$ dans la ligne précédente, qui est juste. L'erreur se
détecte sans calcul : le corrigé affirme dans la phrase suivante que la fonction est
\textit{décroissante} de $\alpha$ et que le maximum est atteint en $\alpha=0{,}4$, or l'expression
écrite est croissante. C'est bien la version corrigée qui est décroissante, et c'est elle qui donne
le résultat encadré $L_1>108~\mu$H, lequel est juste.

\medskip
\erreur{} \textbf{question 4, la seconde moitié n'est pas traitée.} « Quel est l'intérêt du
convertisseur SEPIC ? » reste sans réponse, alors que c'est le cœur de l'exercice : le gain
$\alpha/(1-\alpha)$ vaut $1$ pour $\alpha=1/2$, il est donc \textbf{abaisseur pour $\alpha<1/2$ et
élévateur pour $\alpha>1/2$}, et cela \textbf{sans inversion de signe} contrairement au Buck-Boost.
C'est exactement ce qu'exige un supercondensateur dont la tension traverse les $36$~V en se
déchargeant, et c'est ce qui justifie tout le préambule de l'énoncé.

\medskip
\erreur{} \textbf{question 1, période fausse.} « un cycle de durée $T=1/300\times10^5$~s » : l'exposant
est faux et la parenthèse manquante. Pour $f=300$~kHz, $T=1/f=3{,}3~\mu$s.

\medskip
Deux relâchements de calcul, sans conséquence sur les résultats :
\begin{itemize}
	\item question 2, le facteur $L_1$ disparaît entre $\frac{1}{T}\int_0^T L_1\frac{\dif I_{L_1}}{\dif t}\dif t$
	et $\frac{1}{T}\big(I_{L_1}(T)-I_{L_1}(0)\big)$ ;
	\item question 4, $\langle V_{L_2}\rangle=-\alpha T V_e+(1-\alpha)TV_s$ : il manque le $1/T$. C'est
	l'aire algébrique qui est nulle, pas cette expression qui serait une moyenne.
\end{itemize}

\medskip
\verifie{} : le bilan de puissance et $I_{e,\max}=V_sI_s/V_e=13{,}3$~A, la maille
$V_e=V_{L_1}+V_{C_1}+V_{L_2}$ relue sur le schéma avec les conventions d'orientation des flèches,
d'où $\langle V_{C_1}\rangle=V_e$, le chronogramme ($-V_e$ puis $+V_s$, cohérent avec la figure), le
gain $V_s/V_e=\alpha/(1-\alpha)$, les deux rapports cycliques $\alpha=0{,}4$ et $\alpha=4/7$, la
borne $L_1>108~\mu$H (recalculée par les deux chemins, avec $V_e$ et avec $\alpha$), et
$C_2>46~\mu$F.

\begin{itemize}
	\item \suggestion{} \textbf{La question 7 ne traite pas le cas le plus défavorable.} L'énoncé
	demande de limiter « l'ondulation \textit{maximale} » tout en fixant $V_e=36$~V, ce qui est
	contradictoire. Or $\Delta V_s=\alpha I_s/(fC_2)$ croît avec $\alpha$ : le pire cas est
	$\alpha=4/7$, c'est-à-dire $V_e=27$~V, et il conduit à $C_2>53~\mu$F et non $46~\mu$F. Soit on
	aligne la question sur la question 6 en écrivant « quel que soit $V_e$ », soit on retire le mot
	« maximale ». En l'état un élève rigoureux se trouve piégé.
	\item \suggestion{} Le \textbf{régime permanent périodique} n'est jamais posé en hypothèse, alors
	que toute la question 2 et l'ensemble des bilans reposent dessus. Une ligne dans la liste des
	hypothèses suffirait.
	\item \suggestion{} La convention d'ondulation n'est pas précisée. Le corrigé prend $\Delta i_e$
	crête à crête, ce qui est l'usage, mais l'énoncé ne le dit pas et la valeur change d'un facteur
	deux selon la lecture.
	\item \suggestion{} Notation $I_{e,\max}$ : la question demande la \textit{valeur moyenne} maximale,
	le corrigé répond correctement, mais l'indice « max » évoque plutôt la valeur instantanée crête,
	qui vaut ici $1{,}05\,\langle I_e\rangle$. Un indice du type $\langle I_e\rangle_{\max}$ lèverait
	l'ambiguïté.
	\item \suggestion{} Les figures écrivent $V_S$, $I_S$, $V_E$ et $V_{L2}$ là où le texte écrit
	$V_s$, $I_s$, $V_e$ et $V_{L_2}$. À harmoniser, au moins dans le texte.
	\item \suggestion{} Les deux fichiers de figures ont leurs mots inversés :
	\texttt{conversion\_electronique\_puissance1\_1.png} et
	\texttt{conversion\_puissance\_electronique1\_2.pdf}. Autant renommer le second.
	\item \suggestion{} L'exercice dimensionne $L_1$ et $C_2$ mais laisse de côté $L_2$ et $C_1$. Si
	l'on cherche à l'allonger, c'est l'extension naturelle ; sinon, un mot expliquant pourquoi ces
	deux-là suffisent au dimensionnement demandé.
	\item \suggestion{} Aucune application numérique sur les ondulations elles-mêmes une fois les
	composants choisis, alors que c'est ce qui rend le résultat parlant.
	\item \suggestion{} Coquilles : « converstisseur », « Convertor » (SEPIC = \textit{Single Ended
	Primary Inductor \textbf{Converter}}), « l'intesité », « quelque soit » pour « quel que soit »,
	« A partir » pour « À partir », « L'onudlation », $I_St$ pour $I_st$, et « bobine d'induction »
	là où « bobine » suffit.
	\item \suggestion{} Si l'exercice provient d'un sujet de concours, l'attribuer comme le fait le
	\S~1.1 aiderait à retrouver l'énoncé complet.
\end{itemize}

\section{Magnétostatique}

\traite{} \textbf{le 14 septembre 2026, par réécriture.} Le défaut étant structurel, l'exercice a été
refait plutôt que rapiécé. Il s'appelle désormais « Solénoïde épais à densité de courant variable »
et compte quatre questions : symétries, théorème d'Ampère donnant $B(r)-B(0)$, puis
\textbf{la question qui manquait} (champ nul à l'extérieur, d'où $B(0)$ et $\vec B$ partout), et
enfin la vérification par superposition de solénoïdes emboîtés, qui remplace l'ancienne question 3.
L'énoncé décrit maintenant un bobinage à densité de spires variable et non un conducteur massif, et
porte des valeurs numériques ($a=10$~cm, $j_0=5{,}0$~A$\cdot$mm$^{-2}$, d'où $B(0)=0{,}31$~T).
Deux remarques pour le colleur ont été ajoutées : l'indétermination du cylindre rigoureusement infini,
et la réalisation pratique du dispositif.

\medskip
Les deux figures à main levée ont été refaites, la première en TikZ sous forme de \textbf{coupe
méridienne} (le courant orthoradial y traverse le plan de figure, ce qui est nettement plus lisible
que la perspective d'origine) et la seconde en matplotlib. Les originales restent accessibles dans
l'historique git et sous \texttt{Colle\_compilation/figures/magnetostat\_orthoradial*.pdf}.

\medskip
Le détail de ce qui n'allait pas est conservé ci-dessous.

\exo{4.1 Champ magnétique dans un cylindre parcouru par un courant orthoradial}

Le calcul est juste, et j'ai vérifié la cohérence des signes entre l'orientation du contour $ABCD$ de
la figure et le sens du courant enlacé. Mais l'exercice a un défaut de construction qui dépasse la
correction ligne à ligne, et qui justifie à mon sens une réécriture.

\medskip
\erreur{} \textbf{L'exercice contourne sa propre question centrale.} Le seul pas physique
intéressant est la détermination de la constante d'intégration. Le cylindre étant long mais fini, le
champ tend vers zéro loin de lui ; le théorème d'Ampère montre par ailleurs que $B$ est
\textit{uniforme} pour $r>a$ ; donc $B(r>a)=0$, ce qui fixe
\begin{align*}
	B(0)=\frac{\mu_0j_0a}{2}
	\qquad\text{et}\qquad
	B(r)=\frac{\mu_0j_0}{2a}\left(a^2-r^2\right)
\end{align*}
Or la question 2 demande $B(r)$ « en fonction de la valeur du champ en $r=0$ », donc laisse la
constante indéterminée, et la question 3 introduit un champ extérieur pour imposer $B(a)=0$,
condition qui est \textbf{déjà vraie sans aucun champ extérieur}.

\medskip
\erreur{} \textbf{Conséquence : la question 3 est vide.} Le corrigé répond
$B_{\mathrm{ext}}=-B(0)+\mu_0j_0a/2$, qui n'est pas un résultat : c'est une expression d'un
inconnu en fonction d'un autre inconnu. Si l'on fait la physique correctement, $B(0)=\mu_0j_0a/2$ et
donc $B_{\mathrm{ext}}=0$. La question demande d'ajouter un champ pour obtenir une situation qui est
déjà celle du problème.

\medskip
\erreur{} \textbf{La seconde figure vend la mèche et contredit le texte.} Le tracé placé à la fin de
la question 2 part de $\mu_0j_0a/2$ sur l'axe et s'annule en $r=a$ : c'est le résultat de la
question 3, dessiné comme s'il était celui de la question 2, dont le texte affirme pourtant que
$B(0)$ est indéterminé. Cette figure est la preuve que la confusion est réelle et pas seulement
formelle.

\medskip
\verifie{} : les invariances et le plan de symétrie (le plan $z=\mathrm{cte}$ est bien plan de
symétrie de la distribution, donc $\vec B$, pseudo-vecteur, lui est perpendiculaire, d'où
$\vec B=B(r)\vec e_z$), le courant enlacé $h j_0 r^2/2a$, la circulation $-hB(r)+hB(0)$ avec
l'orientation effectivement dessinée, et les deux expressions de $B(r)$ pour $r<a$ et $r>a$.
Recoupé par une seconde méthode : en voyant le cylindre comme un empilement de solénoïdes
élémentaires de courant surfacique $j(r')\dif r'$, chacun créant $\mu_0 j(r')\dif r'$ à l'intérieur
et rien à l'extérieur, on retrouve $B(r)=\int_r^a\mu_0j_0\frac{r'}{a}\dif r'=\frac{\mu_0j_0}{2a}(a^2-r^2)$.

\begin{itemize}
	\item \suggestion{} \textbf{Le modèle physique n'est pas défendable tel qu'énoncé.} « Un cylindre
	conducteur dans lequel circule une densité volumique de courant orthoradiale » : dans un
	conducteur massif, un courant azimutal permanent supposerait un champ électrique orthoradial, donc
	une circulation non nulle de $\vec E$ sur une boucle fermée, ce qui est impossible en régime
	stationnaire. La distribution est mathématiquement licite ($\diver\vec\jmath=0$), mais sa
	réalisation est un \textbf{bobinage} dont la densité de spires varie avec $r$, pas un bloc
	conducteur. C'est exactement la question que pose un bon élève.
	\item \suggestion{} \textbf{Les deux figures sont dessinées à main levée}, à la tablette, ce qui
	détonne dans un recueil par ailleurs entièrement en Inkscape et matplotlib. La première est peu
	lisible et, surtout, représente le cylindre comme un disque aplati alors que l'énoncé pose
	$L\gg a$. L'annotation de la seconde est indéchiffrable à la taille d'impression et se lit
	volontiers $\propto-1/r^2$, alors que la dépendance est en $-r^2$.
	\item \suggestion{} Aucune valeur numérique. Avec $a=10$~cm et $j_0=5\times10^{6}$~A$\cdot$m$^{-2}$,
	soit $5$~A$\cdot$mm$^{-2}$, ordre de grandeur d'un bobinage bien refroidi, on obtient
	$B(0)=0{,}31$~T : un électroaimant de laboratoire honnête. Cela ancre l'exercice.
	\item \suggestion{} Le profil obtenu mérite d'être commenté : le champ est \textbf{maximal sur
	l'axe et nul au bord}, à l'inverse du fil massif à courant axial que les élèves connaissent. La
	lecture en solénoïdes emboîtés l'explique d'une phrase : en $r$, seules les couches extérieures
	contribuent.
	\item \suggestion{} Coquilles et notations : « l'axe $O_z$ » pour $Oz$, « A l'aide » pour
	« À l'aide », « Quel est l'expression » pour « Quelle », « de sorte à ce que » pour « de sorte
	que », $\oint\dif\vec l\cdot\vec B$ au lieu de $\oint\vec B\cdot\dif\vec l$, les $\dif$ droits, et
	$B_{ext}$ en italique. Le point le plus gênant est l'usage de $B$ à la fois comme sommet du
	contour et comme norme du champ, dans la même ligne de calcul. Enfin l'équation du théorème
	d'Ampère est en \texttt{align} alors que tout le reste est en \texttt{align*}, ce qui produit un
	numéro d'équation isolé.
\end{itemize}

\section{Induction}

\traite{} \textbf{le 14 septembre 2026.} Les cinq erreurs ci-dessous sont corrigées dans
\texttt{main.pdf} : dépendance temporelle rétablie dans $E(r,t)$ et donc facteur $1/2$ de la valeur
moyenne, expression de $j(r,t)$ ajoutée, application numérique refaite, numérotation alignée sur
l'énoncé, notations $\sigma$, $P_{\mathrm{Joule}}$ et $\iint$ harmonisées.

\medskip
Trois changements de fond. \textbf{Une cinquième question a été ajoutée}, qui fait calculer
l'épaisseur de peau et vérifier a posteriori l'hypothèse admise au début : c'est le point le plus
instructif de l'exercice, et il se retourne en argument positif puisque le chauffage par induction
repose précisément sur $\delta\ll e$. \textbf{Le rayon de l'inducteur est passé à $a=7{,}0$~cm}
contre $b=10$~cm pour le disque, de sorte que le terme en $\ln(b/a)$, qui représentait la moitié du
calcul, ne soit plus annulé. \textbf{Le solénoïde a été rallongé} à $\ell=50$~cm, ce qui rend la
formule $B=\mu_0NI/\ell$ défendable, avec une remarque rappelant que seul $nI=B_m/\mu_0$ a un sens
physique et qu'un inducteur réel est une spirale plate.

\medskip
L'objet chauffé est désormais décrit comme un \textbf{disque de cuivre} et non comme un fond de
casserole : $\sigma=6{,}0\times10^{7}$~S$\cdot$m$^{-1}$ est exactement la conductivité du cuivre. Une
remarque explique en quoi les valeurs de l'énoncé ne décrivent pas une plaque de cuisine.

\medskip
\textbf{La version antérieure est conservée} dans
\texttt{relecture/versions/induction\_avant-correction\_2026-09-14.tex} : il suffit de la recopier sur
\texttt{themes/induction.tex} pour revenir en arrière. La figure n'a pas été touchée.

\medskip
Le détail des erreurs est conservé ci-dessous.

\exo{5.1 Chauffage par induction}

Le sujet est excellent et la figure est claire et correcte. Le calcul de la puissance est juste dans
sa structure, y compris la séparation en deux domaines $r<a$ et $r>a$, que j'ai refaite. Mais le
corrigé perd la dépendance temporelle dès la première ligne, et l'application numérique ne suit pas.

\medskip
\erreur{} \textbf{Le facteur $1/2$ de la valeur moyenne manque.} Le corrigé écrit
$E(r)=rB_m\omega/2$, sans la dépendance temporelle. Avec $\vec B=B_m\cos(\omega t)\vec e_z$, la loi de
Faraday donne en réalité
\begin{align*}
	E(r,t)=\frac{rB_m\omega}{2}\sin(\omega t) \qquad (r<a)
\end{align*}
La puissance faisant intervenir $E^2$, sa moyenne temporelle vaut la moitié de la valeur calculée.
La question demandant explicitement $\langle P_{\mathrm{Joule}}\rangle$, il faut
\begin{align*}
	\langle P\rangle=\frac{\pi}{4}\,e\,\sigma\,a^4B_m^2\omega^2\left[\frac14+\ln\frac{b}{a}\right]
\end{align*}
Ce que donne le corrigé est la puissance instantanée maximale.

\medskip
\erreur{} \textbf{La question 2 n'est pas traitée.} Elle demande $j(r,t)$ ; le corrigé calcule $E(r)$
puis enchaîne sur la puissance. L'expression $j(r,t)=\sigma\frac{rB_m\omega}{2}\sin(\omega t)$
n'apparaît jamais, alors que c'est l'objet de la question.

\medskip
\erreur{} \textbf{L'application numérique ne correspond pas à la formule.} Avec $a=b=10$~cm on a
$\ln(b/a)=0$ et le crochet se réduit à $1/4$. En reprenant toutes les données de l'énoncé je trouve
\begin{align*}
	P_{\mathrm{max}}=707~\mathrm{W} \qquad\text{et}\qquad \langle P\rangle=353~\mathrm{W}
\end{align*}
et non les $3769$~W annoncés, que je n'ai pu retrouver par aucune lecture plausible des données. Par
ailleurs la question demande un \textit{ordre de grandeur} et la réponse est donnée à quatre chiffres
significatifs.

\medskip
\erreur{} \textbf{Le « solénoïde » de la dernière question n'en est pas un.} Il est long de $5$~cm et
doit englober la zone de champ de rayon $a=10$~cm : son rapport longueur sur rayon vaut au plus
$0{,}5$, soit exactement l'inverse de la condition du solénoïde long. La formule $B=\mu_0NI/L$ n'y est
pas applicable. Le résultat $I\simeq0{,}04$~A est cohérent avec la formule, mais la formule ne l'est
pas avec la géométrie. Accessoirement, une plaque à induction réelle utilise une spirale plate, pas
un solénoïde.

\medskip
\erreur{} \textbf{Décalage de numérotation.} L'énoncé compte cinq questions, le corrigé quatre items :
l'item 2 répond aux questions 2 et 3 à la fois, et tout est décalé ensuite.

\medskip
Trois relâchements de notation : la conductivité est notée $\sigma$ dans l'énoncé et $\gamma$ dans le
corrigé ; la puissance \textit{totale} est appelée $P_{vol}$, ce qui la confond avec une densité
volumique ; et l'intégrale de flux est écrite avec $\oiint$, réservé aux surfaces fermées, alors que
$\Sigma$ est un disque. S'ajoute une parenthèse non fermée à la question 1.

\medskip
\verifie{} : la forme $\vec\jmath=j(r,t)\vec e_\theta$, les deux expressions de $E(r)$ pour $r<a$ et
$r>a$, les deux intégrales de puissance et leur regroupement en
$\frac{\pi}{2}e\sigma a^4B_m^2\omega^2\left[\frac14+\ln\frac ba\right]$, et le courant
$I=BL/(\mu_0N)=0{,}040$~A.

\begin{itemize}
	\item \suggestion{} \textbf{L'hypothèse centrale est violée par les valeurs numériques choisies.}
	L'énoncé admet que le champ induit est négligeable devant le champ appliqué, c'est-à-dire que le
	champ pénètre uniformément toute l'épaisseur. Or l'épaisseur de peau vaut
	\begin{align*}
		\delta=\sqrt{\frac{2}{\mu_0\sigma\omega}}=0{,}52~\mathrm{mm}
	\end{align*}
	contre $e=3$~mm : le courant est en réalité confiné dans une couche six fois plus mince que
	l'épaisseur supposée traversée. Le corrigé conclut pourtant que le résultat « est cohérent avec la
	puissance maximale des plaques à induction du commerce », ce qui donne une fausse confiance dans un
	modèle invalide aux valeurs retenues. \textbf{C'est une très bonne question à ajouter} : vérifier a
	posteriori la validité de l'hypothèse en calculant $\delta$. C'est d'ailleurs précisément parce que
	$\delta\ll e$ que le chauffage par induction fonctionne.
	\item \suggestion{} Les données ne décrivent pas une plaque à induction. $\sigma=6{,}0\times10^{7}$~S$\cdot$m$^{-1}$
	est la conductivité du \textbf{cuivre} ; un fond de casserole compatible induction est en acier
	ferromagnétique ($\sigma\sim10^{6}$~S$\cdot$m$^{-1}$, $\mu_r$ de l'ordre de $100$). Et
	$B_m\sim10^{-4}$~T est cent fois trop faible : une plaque réelle produit quelques $10^{-2}$~T, ce
	que confirme le résultat de la dernière question, $0{,}04$~A, quand une plaque réelle tire des
	dizaines d'ampères.
	\item \suggestion{} Poser $a=b$ dans l'application numérique annule le terme en $\ln(b/a)$ et rend
	inutile toute la seconde intégrale, celle du domaine $r>a$, qui a pourtant demandé la moitié du
	travail. Prendre par exemple $a=7$~cm et $b=10$~cm ferait travailler le résultat complet.
	\item \suggestion{} L'argument de symétrie de la question 1 est affirmé plutôt que mené. Le plus
	court est d'invoquer l'analogie explicitement : $-\partial\vec B/\partial t$ joue le rôle de
	$\mu_0\vec\jmath$ dans le problème magnétostatique du cylindre parcouru par un courant uniforme
	selon $\vec e_z$, dont le champ est orthoradial. Le corrigé l'esquisse mais parle de « théorème
	d'Ampère » là où il s'agit de la forme intégrale de Maxwell-Faraday.
	\item \suggestion{} Typographie : « A l'aide » pour « À l'aide », $P_{Joule}$ en italique,
	$\vec{e_z}$ au lieu de $\vec{e}_z$, \texttt{rad.s\textasciicircum\{-1\}} et
	\texttt{S.m\textasciicircum\{-1\}} avec des points au lieu de $\cdot$, aucune espace insécable
	avant les unités, et un point final manquant à la question 4 du corrigé.
\end{itemize}

\section{Ferromagnétisme}

\traite{} \textbf{le 14 septembre 2026.} Les six erreurs ci-dessous sont corrigées dans
\texttt{main.pdf}. \textbf{La question 4 est désormais traitée} : matériau dur, condition
$|I|<H_cL/N$, et le contraste chiffré entre un AlNiCo ($|I|<50$~A) et un fer doux ($|I|<10$~mA), qui
est le c\oe ur de la question. \textbf{Les deux appels de figures sont remis dans le bon ordre.} Le
moment magnétique vaut maintenant $\mu_r NIS$ et non $NIS$, est noté $m$ et non $M$, et le champ
dipolaire sur l'axe porte son facteur $2$. L'argument de continuité est reformulé en composante
normale.

\medskip
Deux ajouts. \textbf{Une figure de cycle d'hystérésis} a été tracée et placée à la question 2, avec
le point de fonctionnement à courant négatif : sans elle, la question 4 n'avait aucun support
graphique. Et l'exercice porte désormais des valeurs numériques ($L=20$~cm, $R=1{,}0$~cm, $N=200$,
$I=1{,}0$~A), d'où $H=1{,}0\times10^{3}$~A$\cdot$m$^{-1}$, $B\simeq1{,}3$~T dans le matériau contre
$1{,}3$~mT sans noyau, un moment de $63$~A$\cdot$m$^{2}$ et un champ lointain de $13~\mu$T à $1$~m,
soit le quart du champ terrestre.

\medskip
Un point n'a pas été suivi. La question 4 demandait de \textit{tracer} les trois champs dans le cas
$I<0$ ; le corrigé les décrit précisément par rapport à la figure de la question 3 plutôt que d'ajouter
un cinquième tracé. Une figure matplotlib jurerait avec les trois figures Inkscape existantes : si tu
veux la faire, le mieux est de reprendre \texttt{ferromagnetisme1\_2} en basculant le plateau de $H$
sous l'axe.

\medskip
\textbf{La version antérieure est conservée} dans
\texttt{relecture/versions/ferromagnetisme\_avant-correction\_2026-09-14.tex}.

\medskip
Le détail des erreurs est conservé ci-dessous.

\exo{6.1 Matériau ferromagnétique}

La progression $\vec H$, puis $\vec M$, puis $\vec B$, et enfin l'aimantation rémanente, est une très
bonne idée : c'est exactement l'ordre dans lequel il faut penser un milieu magnétique. Mais le corrigé
s'arrête avant la fin et comporte deux erreurs de fond sur le champ lointain.

\medskip
\erreur{} \textbf{La question 4 n'est pas corrigée.} Le corrigé en recopie l'énoncé mot pour mot, et
s'arrête là. C'est pourtant la seule question qui exploite réellement le ferromagnétisme, et de loin
la plus intéressante. Ce qu'elle appelle : il faut un matériau \textbf{dur}, à cycle d'hystérésis
large. Avec $I<0$ on a $H=NI/L<0$, mais tant que $|H|<H_c$ on reste sur la branche supérieure du
cycle, $M$ reste voisine de l'aimantation rémanente $M_r>0$, et $B=\mu_0(H+M)>0$ puisque
$M\gg|H|$. « Légèrement » signifie donc
\begin{align*}
	|I|<\frac{H_c L}{N}
\end{align*}
Avec un aimant dur (ferrite, AlNiCo, NdFeB), $H_c$ va de $10^{4}$ à $10^{6}$~A$\cdot$m$^{-1}$.

\medskip
\erreur{} \textbf{Les deux figures sont interverties.} \texttt{ferromagnetisme1\_2.pdf}, appelée à la
question 2 qui ne demande que $\vec M$, trace en fait $M$, $H$ et $B/\mu_0$ sur la totalité de l'axe :
c'est la réponse complète de la question 3. Et \texttt{ferromagnetisme1\_3.pdf}, appelée à la
question 3 juste après la phrase « ce qui permet d'avoir l'allure générale de la courbe ci-dessous »
à propos de $\vec B$, ne montre ni $B$ ni le champ extérieur. Il suffit d'échanger les deux appels.

\medskip
\erreur{} \textbf{Le moment magnétique oublie le matériau.} Le corrigé pose $m=NIS$, qui est le moment
du \textit{bobinage seul}. Avec un noyau aimanté, les courants liés ajoutent un courant surfacique
d'intensité linéique $M$, et le moment total vaut
\begin{align*}
	m=(H+M)\,LS=\mu_r\,NIS
\end{align*}
soit environ mille fois plus. C'est précisément l'effet que l'exercice cherche à illustrer, et il
disparaît au moment de calculer le champ lointain.

\medskip
\erreur{} \textbf{Le facteur $2$ du champ dipolaire sur l'axe manque.} Sur l'axe d'un dipôle,
$B=\dfrac{\mu_0}{4\pi}\dfrac{2m}{z^3}$, et non $\dfrac{\mu_0}{4\pi}\dfrac{m}{z^3}$. Avec la correction
précédente, le champ lointain vaut donc $B=\dfrac{\mu_0\mu_r NIR^2}{2z^3}$.

\medskip
\erreur{} \textbf{Le moment magnétique est noté $M$}, c'est-à-dire du même symbole que l'aimantation,
dans le paragraphe qui manipule les deux.

\medskip
\erreur{} \textbf{« le champ magnétique doit être nécessairement continu (à cause de $\diver\vec B=0$) ».}
L'équation $\diver\vec B=0$ n'impose la continuité que de la composante \textit{normale}. La composante
tangentielle de $\vec B$ est d'ailleurs discontinue sur la surface latérale du solénoïde, là où
circulent les courants. La conclusion est juste sur l'axe, où il s'agit bien de la composante normale
aux faces d'extrémité, mais l'argument tel qu'il est écrit est trop fort.

\medskip
\verifie{} : $H=NI/L$, l'argument de symétrie (le plan $z=\mathrm{cte}$ contient les courants, donc
$\vec H$, pseudo-vecteur, lui est perpendiculaire), la relation $\vec B=\mu_0(\vec H+\vec M)$,
l'approximation $\vec B\simeq\mu_0\vec M$ dans le matériau, la décroissance en $1/z^3$, et le contenu
physique de la figure \texttt{1\_2} : $B/\mu_0$ continu, $H$ et $M$ discontinus en sens opposés aux
deux extrémités.

\begin{itemize}
	\item \suggestion{} \textbf{La question 2 répond en matériau linéaire alors que la question 4 exige
	l'hystérésis.} Le corrigé invoque un « matériau ferro. linéaire et doux » de perméabilité relative
	$10^{3}$, c'est-à-dire précisément le contraire de ce dont la question 4 a besoin : aimantation
	rémanente et champ coercitif. Il faut introduire dès la question 2 la courbe de première
	aimantation, la saturation et le cycle d'hystérésis, sans quoi la dernière question n'a aucun
	support.
	\item \suggestion{} La question 3 demande explicitement « l'allure de $\vec H$ sur la totalité de
	l'axe », mais aucune figure ne trace $H$ seul. Le corrigé le décrit en mots, ce qui est dommage
	pour une question de tracé.
	\item \suggestion{} La question 1 utilise l'idéalisation du solénoïde infini ($H$ nul à
	l'extérieur), et la question 3 étudie justement le champ à l'extérieur d'un solénoïde fini. La
	tension mérite une phrase : on travaille d'abord dans la limite infinie, puis on revient au
	solénoïde réel.
	\item \suggestion{} Aucun résultat encadré, aucune valeur numérique, aucune remarque pour le
	colleur, contrairement au standard du recueil. Un ordre de grandeur sur $M_r$ et $H_c$ rendrait la
	question 4 concrète.
	\item \suggestion{} L'origine de $z$ n'est pas marquée sur la figure du montage, alors que les deux
	tracés placent le solénoïde entre $z=0$ et $z=L$.
	\item \suggestion{} Coquilles : « connaitre », « un des côté de longueur $l$ en confondu avec
	l'axe » (deux fautes dans la même proposition), « de sorte à ce que », « décroit », l'abréviation
	« ferro. », $\vec{e_z}$ au lieu de $\vec{e}_z$, « En déduire l'allure $\vec H$ » où il manque
	« de », et l'ordre des mots de la phrase d'introduction, difficile à lire.
\end{itemize}

\end{document}
